\begin{exercise}[The Sagnac time delay]{McGreevy, Problem Set 4}
Light is sent in both directions around a circular ring of radius \(R\) that
rotates with angular velocity \(\Omega\).  Compute the return-time difference
measured by a clock on the ring, first exactly in the inertial frame and then
to leading order in \(\Omega R\).  Explain why the effect cannot be removed by
a single rotating inertial frame.
\end{exercise}

\begin{exercise}[Fermi--Walker transport]{Cambridge Part III, General Relativity examples}
Derive the transport law for a torque-free spin \(S^\mu\) carried by an
accelerated observer with velocity \(u^\mu\) and acceleration \(a^\mu\),
requiring preservation of \(S\cdot u=0\) and absence of spatial rotation in
the instantaneous rest frame.
\end{exercise}

\begin{exercise}[How large may Einstein's elevator be?]{MIT 8.962, Problem Set 4}
Two neighboring freely falling particles have separation \(\xi^i\).  Use
geodesic deviation to estimate the fractional difference of their
accelerations across an elevator of size \(L\) near a spherical body of mass
\(M\) at radius \(r\).  State the condition under which a uniform
accelerating laboratory is a good approximation.
\end{exercise}

\begin{exercise}[Exact redshift in a Rindler laboratory]{Tong, Cambridge Part III Sheet 2}
For
\[
\dd s^2=-(1+a\xi)^2\dd\eta^2+\dd\xi^2+\dd y^2+\dd z^2,
\]
compare the frequencies of a photon measured by static observers at
\(\xi_1\) and \(\xi_2\).  Recover the weak-field gravitational-redshift
formula and locate the acceleration horizon.
\end{exercise}

\begin{exercise}[Composition dependence and the E\"otv\"os parameter]{Cambridge Part III, General Relativity examples}
Suppose two test bodies have gravitational-to-inertial mass ratios
\(1+\delta_A\) and \(1+\delta_B\).  Derive their E\"otv\"os parameter
\(\eta_{AB}=2|a_A-a_B|/(a_A+a_B)\) to first order.  Explain which form of the
equivalence principle such an experiment tests.
\end{exercise}

\begin{exercise}[The equivalence-principle estimate of light bending]{MIT 8.962, Problem Set 4}
Model a light ray passing a mass \(M\) with impact parameter \(b\) as a
particle moving at unit speed through the Newtonian transverse acceleration.
Compute the deflection.  Compare it with the general-relativistic result and
identify the missing physical effect.
\end{exercise}

\begin{exercise}[Born-rigid acceleration]{McGreevy, Problem Set 3}
A rod accelerates along its length while maintaining constant proper distance
between neighboring points.  Show that its worldlines are curves of constant
\(\xi\) in Rindler coordinates and that the proper acceleration must vary as
\(a(\xi)=a_0/(1+a_0\xi)\).  Why can the whole rod not have the same proper
acceleration?
\end{exercise}

\begin{exercise}[Clock comparison near Earth]{Cambridge Part III, General Relativity examples}
In the weak terrestrial field, compare clocks separated by a small height
\(h\).  Derive the accumulated proper-time difference after coordinate time
\(T\), and estimate its order of magnitude for \(h=1\,\mathrm{m}\) and
\(T=1\) day.
\end{exercise}

\begin{exercise}[Inertial forces in a rotating frame]{McGreevy, Problem Set 4}
Transform Minkowski spacetime to cylindrical coordinates rotating with
angular velocity \(\Omega\).  Read off the proper time of a stationary
rotating observer and derive the nonrelativistic Coriolis and centrifugal
accelerations from the geodesic equation.
\end{exercise}

\begin{exercise}[A gravitational compass]{Cambridge Part III, 2018 exam}
In a local freely falling orthonormal frame, a family of nearby test masses is
released from rest with several independent separation vectors.  Show how
their relative accelerations determine the symmetric tidal tensor
\(E_{ij}=R_{0i0j}\).  How many independent release directions are needed in
vacuum?
\end{exercise}
